\section{总结与展望}

\subsection{工作总结}

本文完成了一套图像级深度伪造检测的完整流程：按身份隔离的数据管线、RGB 基线、双流与门控、多维度评测、可解释性验证和可演示系统。主要结论是：

\begin{enumerate}
  \item 数据量对同库性能影响显著（280$\to$580 身份，AUC $+0.021$，TPR@1\% $+23$ pt）；
  \item 频域分支与质量门控在本数据上没有稳定收益，门控退化为常数；
  \item 真假频谱差异主要来自压缩块效应，而非生成器指纹；
  \item 未见伪造方法上模型接近随机，存在明显的捷径学习；
  \item 系统在单图 11 ms 下可完成检测、解释与报告生成。
\end{enumerate}

\subsection{局限性}

\begin{itemize}
  \item 数据仅覆盖 FF++ C23 的 5 种方法，未包含扩散生成与商用工具；
  \item 图像级检测，未利用视频时序与音频；
  \item 门控设计过于简单，未验证多维质量描述子；
  \item 跨库与开集泛化仍是短板；
  \item 未做多随机种子重复，部分差异可能落在噪声范围内。
\end{itemize}

\subsection{未来工作}

\begin{itemize}
  \item 引入 SBI 等自混合增广，提升对未见方法的泛化；
  \item 测试时 BN 适应 / 域适应，缓解跨库偏移；
  \item 用多维质量描述子（高频能量 + 噪声估计 + 模糊度）重做门控；
  \item 扩展到视频时序与多模态；
  \item 结合重构误差（DIRE/AEROBLADE）做通用 AIGC 检测。
\end{itemize}
